% ch01_intro.tex: Chapter 1, Introduction.
% Every number is traced in ch01_intro_NOTES.md. Framing per OUTLINE s7 risk 1: a benchmark plus one confirmed recipe gain.
\chapter{Introduction}\label{ch:intro}

\section{Electroluminescence inspection and the task}\label{sec:intro-task}

Electroluminescence (EL) images show defects in solar cells that are invisible to the eye, and they are used both before a module
ships and over its life \citep{pratt2023benchmark}. A module driven in forward bias emits light where its silicon is connected and
of good quality; cracks appear as dark lines, and a part of a cell that a crack cuts off from the current no longer contributes
power and appears as a dark region, the inactive area \citep{haase2018crackwidth,kontges2011risk}.

This thesis studies pixel-level semantic segmentation of EL images of crystalline silicon cells: each image is grayscale, 512
$\times$ 512~px, and shows one cell plus parts of its neighbours, and each pixel receives one of 24 classes, 12 structural and 12
defects. The main labelled data are the BenchmarkEL release \texttt{20211104} \citep{pratt2023benchmark}, 582 cells split 478 /
54 / 50 into train, validation and test, from five sources with different contrast and labelling conventions. Crack and inactive
are the focus: crack because it occurs in many cells (206 of the 478 training cells under our label set) and is hard to label,
inactive because of all classes it is the most directly tied to lost power. The work was done with an industry partner, a PV recycling and
inspection company; its EL rig was still 2 to 3 months away as of 4 November 2026, so every model here is built and tested on public
data.

The two focus classes are hard in opposite ways (Chapter~\ref{ch:background}). A physical crack is micrometres wide; the camera
shows a dark line with a median full width at half maximum of 4.0~px; the annotators traced a line of median width 1.75~px somewhere on
it, and half of the label points sit more than 1~px off the darkest ridge. A perfect crack mask shifted by 1~px therefore scores
crack IoU 0.559 on the legacy test masks. Inactive regions are blobby and consistently labelled but rare: inactive occurs in 42 of
the 478 training cells.

\section{The goal and what was reached}\label{sec:intro-goal}

The goal set at the start of this work was, in its original wording, ``highest scores, a model which captures almost all of the
cracks, inactive and other important defects balancing the miou for other classes as well''. We made it operational for cracks as
a relaxed recall of at least 0.90 at a relaxed precision of at least 0.90, both within 2~px (relR@2 and relP@2;
Chapter~\ref{ch:background}). That target is not reached. On the sealed set of unseen modules, the final recipe, averaged over
seeds 5 to 9, reaches relP@2 0.5995 and relR@2 0.4157 at argmax; even in distribution, on the cross-validated pool and with the
recipe chosen on the same cells, it reaches 0.9216 and 0.8211. Inactive gains are not detected, and several of the other classes stay weak.

The first obstacle was not a model but a measurement. To know whether a change helps, one needs a number that means the same thing
before and after the change and can be set beside published numbers. For EL segmentation no such number existed, and the aim of
this work widened from a better model to setting ``a benchmark for this very unorganised field''. This thesis therefore reports a
benchmark first, and one model gain measured under it.

\section{Why published EL segmentation numbers do not compare}\label{sec:intro-why}

We read 22 papers that report pixel-level EL segmentation numbers in full text (Chapter~\ref{ch:related}). Their numbers cannot be
compared with each other or with ours, for six reasons that we measured or counted.
\begin{itemize}
  \item \textbf{Releases.} The main dataset exists in three releases (582, 642 and 695 cells) with reshuffled splits: 186 of the 582
  cells change split between the first two, and the test sets of the later releases each hold 51 cells of our training pool.
  \item \textbf{Leakage.} The legacy test set (test50) is not independent of training: 8 of its 50 cells are the same physical
  cell as a pool cell (7 confirmed by pixels, 1 by the cell-table link only), 7 more are neighbouring crops of pool images, and 19
  more share a module with pool cells.
  \item \textbf{Classes and averaging.} Class lists differ (one dataset's ``crack'' is a region class close to crack plus inactive),
  and so does averaging. On test50 (in-distribution, known leakage, context only), the same predictions of our baseline give a
  per-image mean of defect medians of 0.202, 0.405 or 0.541, depending only on how cells without the class are counted; on what are
  presumably the same 50 cells, published crack values are 0.25 \citep{pratt2023benchmark} and 0.58 \citep{jha2025pvs3}, under
  different metrics.
  \item \textbf{Line width.} Strict IoU on 1 to 2~px crack labels measures the annotators' line-width convention: the same baseline
  predictions on our pool score strict crack IoU 0.233, 0.365 and 0.414 against 1-px centreline, v1.1 and 3-px labels, while the
  2~px tolerance metric relF1@2 reads 0.824, 0.846 and 0.849.
  \item \textbf{Label releases.} The same models score strict crack IoU 0.320 on the 2021 labels and 0.365 on the corrected labels
  we use, and relF1@2 0.822 and 0.846.
  \item \textbf{No uncertainty.} 18 of the 22 papers give no interval, seed count or repeat on any number; none states a
  module-disjoint split, uses a distance-tolerant crack metric or reports a number on unseen modules; 2 state release, split and
  metric together.
\end{itemize}
Our reading (judgement): a comparable number must state seven things together, namely dataset release, label set, split, class
list, metric convention, operating-point rule and seed count, and no published EL number we found states all seven.

\section{Approach}\label{sec:intro-approach}

The work began as a study of semi-supervised learning on the legacy split; under the rule of that time, pseudo-labels helped only
at 5\% and 10\% labels, and the framing was dropped. The measurement failures of that period (a test set scored in every run and overlapping
the training pool, effects near the noise, inactive decided by a few cells) motivated the protocol (Chapter~\ref{ch:legacy}). We
then built the instrument before judging methods with it. A dataset audit measured the faults that move scores: the line-width
convention, label errors, the image window, the label release and a module leak across the folds of our earlier split. A protocol
then fixed the seven elements, with five rules that follow from the audit: a cross-validation split of the 532 pool cells grouped
by physical module; crack relF1@2 (tolerance 2~px) as the deciding metric, with strict IoU printed only beside its convention band;
guards against gains bought by wider cracks or more false alarms; margins from an A/A contrast; and sealed test sets opened once,
under a rule registered before any prediction. Under that protocol we ran a suite of published models and recent methods, chose a final recipe on pool cells, and
tested it once on E1, 106 cells of unseen modules (supported, not proven) from a later release. A few-shot track then asked what
labelled cells of E1's layouts add, and a deployment chapter measured cost, severity output and licences.

\section{Contributions}\label{sec:intro-contrib}

\begin{enumerate}
  \item \textbf{A benchmark protocol and its release} (Chapter~\ref{ch:protocol}). The split cv5v3 groups the 532 pool cells by
  physical module into five folds over 285 groups, none of which spans folds; it replaces a split in which 24 modules (91 cells)
  spanned two folds. Crack relF1@2 decides, with a tolerance of 2~px, the 95th percentile of the distance between two annotations of one crack;
  strict crack IoU is printed with a convention band of $\pm$0.019 and decides nothing; the margin is $M=0.010$; every entry fills
  an info sheet with its entry class, cost and seed count. The benchmark is prepared for release as lists and code only, without images, masks
  or weights.
  \item \textbf{A dataset audit} (Chapter~\ref{ch:data}). The audit shows that strict crack IoU measures the line-width convention,
  that 57\% of crack pixels and 71\% of inactive pixels lie outside the centre cell of each image, that crack label errors exist but
  are a minority of the image-supported candidates (corrected share 0.244 [0.105, 0.436]) and move no contrast by more than 0.003,
  that relF1@2 contrasts move by at most 0.0031 under the label variants tried, and that the label release moves the baseline's
  strict crack IoU by +0.045, more than any method gain over it on the pool.
  \item \textbf{A method suite under one protocol} (Chapter~\ref{ch:methods}). Under equal training, two models of the EL papers sit
  below the baseline on crack relF1@2 ($-$0.0773 [$-$0.0892, $-$0.0594] and $-$0.0932 [$-$0.1113, $-$0.0774], 2-seed screens).
  Of the three trained arms read at five seeds (HRNet-W18, weight averaging and their combination), none was confirmed under its
  registered rule; D4 test-time augmentation on the baseline is confirmed on the pool (+0.0134 [+0.0066, +0.0207]), for known
  modules only.
  \item \textbf{One confirmed recipe gain on a sealed set} (Chapter~\ref{ch:e1_result}). On E1, the final recipe C2 (HRNet-W18
  U-Net++ with SWAD weight averaging and D4 test-time augmentation) beats the encoder-preserving baseline on crack relF1@2 by +0.0430
  [+0.0013, +0.1043], CONFIRMED under a rule registered before any prediction, with all guards passing. Its scope travels with it:
  the recipe averaged over seeds 5 to 9, not one checkpoint; argmax only; 101 of the 106 cells from one source (CSIR); the full
  cells carry it (+0.1464 [+0.0420, +0.2278]), while on the half-cut cells it is not detected (+0.0110 [$-$0.0289, +0.0722]). The
  lower end is near 0 and the smallest detectable effect of the read is about 0.074. Most of the gain is the model (+0.0451
  [+0.0173, +0.0699] with the same eight views on both sides); the D4 part is not resolved on E1. C2's readout and views were
  chosen on pool cells, and E1 had seven earlier registered openings for other models.
  \item \textbf{Cost and severity} (Chapter~\ref{ch:deployment}). C2 costs 1,744 GFLOPs per cell, $11.8\times$ the baseline, and
  about 0.17~s of network time per cell on one GPU at batch 1; its five-seed ensemble (8,722 GFLOPs per cell) adds nothing
  measurable over its members on E1 ($-$0.0025 [$-$0.0254, +0.0157]; OP-A, recipe means over seeds 5 to 9; 101 of 106 cells
  CSIR). The severity output (inactive area fraction per cell) is validated
  only for the baseline on pool cells, and the claim that predicted inactive area predicts power loss better than crack length is
  not supported.
  \item \textbf{A few-shot track} (Chapter~\ref{ch:t3b}). A registered out-of-fold test inside E1 asks how much labelled cells of
  its layouts add to the same recipe. It is CONFIRMED: out of fold over 5 folds of E1 (seeds 5 to 7, OP-A, 106 cells, 101 of them
  CSIR), the recipe trained with the labelled E1 cells of the other folds beats the same recipe zero-shot on crack relF1@2 by
  +0.2213 [+0.0977, +0.3504], with both guards passing. The half-cut cells carry it (+0.2530 [+0.0878, +0.3885]), but not those
  without scuff ($-$0.0701 [$-$0.3019, +0.0517]), and half of the labelled images give about the same gain (ALL $-$ HALF +0.0171
  [$-$0.0246, +0.0625], not separated). It shows that labelled cells of a layout help, not that the model handles a new layout
  without them.
\end{enumerate}

\section{What this thesis does not claim}\label{sec:intro-not}

The headline is a benchmark plus one confirmed zero-shot recipe gain, not a model that finds almost all cracks: on E1 the final recipe
finds about 0.42 of the labelled crack pixels within 2~px (relR@2 0.4157, mean over seeds 5 to 9, argmax). We do not claim to beat
any published model; test50 overlaps our training pool, and its numbers are in-distribution, with known leakage, and context
only. We do not claim an inactive gain, a gain for one deployable checkpoint, a zero-shot gain on half-cut cells, or a solution for
the weak classes. The few-shot gain needs labelled cells of the layout and is measured inside E1 only. ``Unseen modules'' is
supported, not proven. From the few-shot opening on, E1 is training data, so the zero-shot
result of Chapter~\ref{ch:e1_result} is the last on these cells. Chapter~\ref{ch:limitations} lists every limit.

\section{Thesis map}\label{sec:intro-map}

Chapter~\ref{ch:background} explains EL imaging, the 24 classes and why crack and inactive labels behave so differently, and
defines the metrics. Chapter~\ref{ch:related} reviews the public EL data and published numbers and shows why they do not compare.
Chapter~\ref{ch:data} describes the data, the label sets and the dataset audit, whose findings are the reasons for the protocol.
Chapter~\ref{ch:legacy} summarises the work done before the protocol, whose measurement failures motivated it.
Chapter~\ref{ch:protocol} states the benchmark protocol. Chapter~\ref{ch:methods} reports the baselines and methods run under it
and the choice of the final recipe. Chapter~\ref{ch:e1_result} gives the sealed E1 result and the public test-set table.
Chapter~\ref{ch:t3b} reports the few-shot track. Chapter~\ref{ch:deployment} covers severity, inference cost and licences.
Chapter~\ref{ch:limitations} collects the limitations, and Chapter~\ref{ch:conclusion} reads the claims table and names future
work. Appendix~\ref{app:tables} holds the tables, Appendix~\ref{app:related} the related-work table, and
Appendix~\ref{app:t21} a note on reading the pre-protocol table.
